\documentclass[a4paper,11pt]{article}

\usepackage{revy}
\usepackage[utf8]{inputenc}
\usepackage[T1]{fontenc}
\usepackage[danish]{babel}
\usepackage{graphicx}
\usepackage{amssymb}


\revyname{MatematikRevyen}
\revyyear{2026}
% HUSK AT OPDATERE VERSIONSNUMMER
\version{0.1}
\eta{$5$ minutter}
\status{Færdig}

\title{Morgenmad på Hilberts Hotel (lang version)}
\author{Julius Pallesen '24}

\begin{document}
\maketitle

\begin{roles}
\role{M} Matematikkandidat, som er ansat på hotellet
\role{H} Hotelejer Hilbert
\role{K} Kok
\role{G} Venlig gæst på værelse $10^{9672}$
\end{roles}

\begin{props}
    \prop{Meget lang brød pose}
    \prop{Ovn}
    \prop{Køleskab}
    \prop{Morgenmad} Rundstykker, juice og 
    \prop{Spiseborde}
    \prop{Buffetbord}
    \prop{Morgenmadsrecepcion}
\end{props}


\scene{Note: Dette er tredje sketch i en serie af tre. Tbh havde jeg lidt svært ved at afslutte trilogien uden den sidste sketch blev for lang. Anyway, jeg har uploadet både en lang og en kort version. Jeg har også en mellemversion liggende, men den må I spørge mig om hvis I vil have den (se note nederst).}
\begin{sketch}
\scene{Vi er restaurenten på Hilberts Hotel. Der er et buffetbord i midten, og et par spiseborde hvor der sidder nogle hotelgæster. Scenen starter in medias res, med at H kommer ind på scenen og M følger efter.}

\says{M} Jamen chef, jeg er slet ikke færdig med at gøre rent på alle værelserne!

\says{H} Nej, men nu har du arbejdet her i et år, og du er ikke engang nået til værelse én oktotrigintimillillion\footnote{En oktotrigintimillillion $= 10^{228000}$ er lidt over det antal værelser man kan nå med ét års harmonisk rengøring (det mere præcise antal er $351$ septentrigintimillinovemcentinovemtrigintilliarder $\approx3,5\cdot10^{227639}$)}. Der er også andre ting der skal gøres her på hotellet. Nu vil jeg gerne have, at du sætter op til morgenmad. Så hvis du lige gider gå ud i køkkenet og gøre hvad end kokken beder dig om. Imens har jeg et hotel at passe.

\scene{H går. M  går ud i køkkenet, hvor K er i gang med at lave ét røræg.}

\says{M} Hej, jeg...

\says{K} Du kan bare tage rundstykkerne ud af ovnen, de er klar. Der skulle gerne være nok til at alle kan få én hver. Eller to. Eller prik prik prik.

\says{M} Men hvordan er det meningen jeg skal flytte uendeligt mange brød.

\says{K} \act{Suk} Det er dig, vi har hørt om, der har været igang med at gøre værelser rent i et helt år? Prøv at høre her, fister. Her på hotellet er det vigtigt at kunne arbejde effektivt. Ellers holder man ikke i branchen. Nu tager du denne her bundløse brødpose, og så hælder du hele bagepladen ned i, og så går du ind med hele posen, dvs. ALLE brødene PÅ ÉN GANG! Tror du, du kan finde ud af det?

\says{M} Jeg kan prøve.

\scene{K rækker M en brødpose der er så lang, at den når ud sidetæppet. M går over til ovnen som vender væk fra publikum, så man ikke helt kan se, hvad der foregår. Posen bliver fyldt med brød, og M slæber med noget besvær $\infty$ kg ind til buffetbordet, og hælder en god sjat ud på et fad og kommer tilbage ud i køkkenet.}

\says{K} Så må du gerne gå ind i køleskabet og tage én flaske appelsinjuice, én flaske æblejuice og én karton yoghurt. Er opgaven forstået?

\says{M} Kun én af hver? Men skal der ikke være til alle?

\says{K} Prøv at høre her, knejt. Har du ikke en kandidatgrad i matematik?

\says{M} Jo, men...

\says{K} Godt. Hvor mange stykker yoghurt er der så i en karton yoghurt?

\says{M} Stykker yoghurt?

\scene{K åbner en yoghurtkartonen.}

\says{K} Her, prøv at tæl.

\says{M} Man kan da ikke tælle yoghurt. Men der er vel en liter?

\says{K} Hør her, mester. Målteori rager mig en papand. Pointen er at der er overtælleligt meget yoghurt til tælleligt mange gæster. Det burde være rigeligt.

\says{M} Øh okay..? Men bliver gæsterne så mætte hvis de skal dele én liter?

\says{K} Som sagt siger de der volumenmål mig ikke en skid. En liter kan blive til ligeså mange liter du har lyst til, hvis bare du tør udvælge de rigtige stykker. Her, tag det her røræg ind. Og så kan du lige tænke lidt over hvor mange penge hotellet sparer ved at lave røræg i stedet for blødkogte æg. Kan du ikke også lige tage imod gæsterne?

\scene{M er forvirret, men går ind med juice, yoghurt og røræg til buffeten og folk går op og tager. M stiller sig ved morgenmadsreceptionen. G kommer ind.}

\says{M} Værelsesnummer?

\says{G} Én milliseksticentiduodecillion. Lige ud.

\says{M} Hold da op! Det er et af de små tal! Du tager bare plads.

\scene{G går over til nærmeste stol. Personen der sidder sig rejser sig og går hen til næste stol, hvor personen der sidder dér rejser sig osv. indtil en går ud af sidetæppet, hvor der skal forestille at være flere borde.}

\scene{G vender sig og snakker til M.}

\says{G} Hey! Er det egentlig ikke dig, alle taler om, som har gjort rent et helt år? Hvad hedder du?

\says{M} Øh jo.. Jeg hedder M.

\says{G} M, fedt at møde dig. Jeg hedder G. Kan du reglerne til whist?

\scene{G tager et kortspil frem og begynder at dele i fire bunker.}

\says{M} Jeg elsker whist! Men jeg må ikke spille det her på hotellet for min chef.

\says{G} Årh. Hold nu op! Nu har du gjort rent et helt år. Det lyder ikke som om du skylder ham chefen noget. Du har fortjent lidt pause. Et lille spil whist, kan da ikke skade. \act{Ud til hele spisesalen:} Er der andre, der vil være med til whist?

\scene{Alle vender sig mod G og siger ja!}

\says{G} Vent lige, går det her overhovedet op?

\says{M} Nej det gør det jo ikke. $\aleph_0+1$ er jo både kongruent med $0$ og $1$ modulo $4$. Det er en modstrid. Det er derfor man ikke kan spille whist her på hotellet.

\says{G} Når ja. Men lad os prøve alligevel. Det kan bare være os fire \act{peger på de fire nærmeste}. Og så kan I fire spille, og så videre. Altså hvad tror du, der sker ved, at der opstår en lille modstrid. At hele hotellet styrter sammen? Hahaha!

\says{M} Altså nu findes det her hotel jo kun i den abstrakte matematiske verden, hvor en modstrid kan ødelægge alt.

\says{G} Jaja whatever.. \act{kigger på sine kort} 9 vip.

\scene{Lys ned efterfulgt af strobelys og en lyd af en kæmpe bygning der styrter sammen og folk der skriger.}

\scene{Lys helt ned.}

\says{AV} Og sådan sluttede tragedien om Hilberts Hotel. Modbevist ved Whist.
\end{sketch}

\scene{Note 1: I den mellemlange version, er der ingen kok, men i stedet er det G der forklarer K det med yoghurten, og det er noget mere kortfattet.}
\scene{Note 2: Da Whist bliver introduceret som Chekhovs pistol i akt 1, er det imperativt at det bliver brugt i akt 3. Hvis denne sketch bliver valgt enkeltvist, ville jeg fjerne Whist-plottet og erstatte slutningen med en god punchline.}


\end{document}
